% Generated draft: requires mathematical and source review.
\documentclass[11pt]{article}
\usepackage[margin=25mm]{geometry}
\usepackage{fontspec}
\setmainfont{Latin Modern Roman}
\usepackage{amsmath,amssymb,mathtools,unicode-math}
\setmathfont{Latin Modern Math}
\usepackage{enumitem,needspace,xurl,hyperref}
\setlength{\parindent}{0pt}
\setlength{\parskip}{5pt}
\setlength{\emergencystretch}{3em}
\newcommand{\N}{\mathbb{N}}
\newcommand{\Z}{\mathbb{Z}}
\newcommand{\Q}{\mathbb{Q}}
\newcommand{\R}{\mathbb{R}}
\newcommand{\C}{\mathbb{C}}
\newcommand{\F}{\mathbb{F}}
\newcommand{\A}{\mathbb{A}}
\newcommand{\PP}{\mathbb{P}}
\newcommand{\E}{\mathbb{E}}
\newcommand{\eps}{\varepsilon}
\newcommand{\dd}{\,\mathrm{d}}
\begin{document}
\section*{UnsolvedMath ID: 9400126}
\textbf{Source problem code:} catalogue entry
\subsection{Definitions and mathematical statement}
For a positive integer $n$, the Mersenne number is $M_n=2^n-1$. Its prime-divisor set is
\[
 D_n=\{\ell:\ell\text{ prime and }\ell\mid2^n-1\}.
\]
The source identifies Gillies' conjecture as concerning the distribution of these prime divisors, but gives no distribution law, asymptotic formula, range, error term, or exponent restriction. It does not say whether $n$ ranges over all integers or only primes, or whether the counted data are divisor counts, factor sizes, or another statistic.

Accordingly the record names a conjecture but does not determine its quantitative mathematical statement. It must not be replaced by the different question whether infinitely many $M_n$ are prime.

The supplied summary also calls infinitude of primes dividing some Mersenne number open. That background claim is erroneous in the unrestricted-index meaning: every odd prime $\ell$ divides $2^{\ell-1}-1$, so such prime divisors already include all odd primes. This observation identifies a source-status error, not a resolution of Gillies' unspecified distribution law. No precise Gillies formula is inferred from the mistaken paraphrase.
\subsection{Short English statement}
State and investigate Gillies' proposed distribution of prime divisors of Mersenne numbers. The actual distribution statement is missing. The summary's replacement by infinitude of all Mersenne divisors is incorrect, and infinitude of Mersenne primes is a separate question. The OPEN-TRIAGE label therefore does not verify the named conjecture's exact current status.
\subsection{Sources}
\begin{itemize}
\item Local UnsolvedMath registry, permanent ID 9400126, source code catalogue entry, ``Gillies' conjecture.'' Dataset curation and metadata: Ulam AI and UnsolvedMath Contributors, CC BY 4.0.
\item Wikipedia contributors, \emph{List of unsolved problems in mathematics}, Number theory, Prime numbers bullet 126, wikitext line 160, pinned revision 1366636767, \url{https://en.wikipedia.org/w/index.php?title=List_of_unsolved_problems_in_mathematics&oldid=1366636767#Number_theory}; MediaWiki-source extraction, accessed 2026-07-29.
\item The summary points to ``Mersenne prime'' on Wikipedia without a primary Gillies reference, exact formula, or linked-page revision.
\item Source provenance stages the candidate pending semantic deduplication and primary-source status review; current status and redistribution terms NEEDS\_REVIEW, separately from dataset attribution.
\end{itemize}
\end{document}
